\part{Using Git with others}

\section{Using Git with others}

\begin{frame}
\frametitle{Working with others}
\begin{itemize}
    \item a networked world
    \item version control becomes more interesting when one works with
        others
    \item author and committer can be different people (usually aren't)
    \item work more often with branches (very powerful feature)
\end{itemize}
\end{frame}

\begin{frame}
\frametitle{Working with others (also online)}

\begin{description}
    \item[\code{git clone}] Clones a pre-existing repository; basically copy
        into a new directory and get the current branch
    \item[\code{git pull}] Pull changes from another repo into the local repo
    \item[\code{git push}] Push changes from a local repo into another
    \item[\code{git format-patch}] Generate a patch file, which one can
        then send as an email attachment
\end{description}
\end{frame}

\begin{frame}
\frametitle{Tagging commits}
\end{frame}

\begin{frame}
\frametitle{Add remote repos to model}
\begin{itemize}
    \item working directory
    \item staging area
    \item local repo
    \item remote repo(s)
\end{itemize}
\end{frame}

\begin{frame}
\frametitle{Cloning an existing repo}
\begin{itemize}
    \item git clone
    \item changes to git config --list
    \item explain meaning of origin
    \item master, origin/master, HEAD
\end{itemize}
\end{frame}

\begin{frame}
\frametitle{Exercise!}
\end{frame}

\begin{frame}
\frametitle{Dealing with conflicts}
\begin{itemize}
    \item can happen when working on same part of a file
    \item happens when working with others
    \item happens when working with branches and stashes (hence can also
        happen when one is working alone!)
\end{itemize}
\end{frame}

\begin{frame}
\frametitle{Exercise!}
\end{frame}

\begin{frame}
\frametitle{Branches}
\begin{itemize}
    \item what are they?
    \item why use them?
    \item localise work on a focussed topic
    \item don't interfere with main development on master
    \item allows devs to work on parts of projects independently
    \item allows for code reviews before merging
\end{itemize}
\end{frame}

\begin{frame}
\frametitle{The default branch: master}

Discuss push/pull just from master?  Then have push/pull concept ready
for quick fix example below
\end{frame}

\begin{frame}
\frametitle{Branches}
\begin{itemize}
    \item git branch
    \item git branch -a
    \item git branch <branchname>
    \item git checkout <branchname>
    \item git checkout -b <branchname>
    \item naming branches -> clarity very helpful
\end{itemize}
\end{frame}

\begin{frame}
\frametitle{Simple branch}
\begin{center}
    \resizebox{!}{0.7\textheight}{%
        \input{figures/simple_branch_tikz.tex}
    }
\end{center}
\end{frame}

\begin{frame}
\frametitle{Examples and exercises!}
\end{frame}

\begin{frame}
\frametitle{Merging branches}
\begin{itemize}
    \item merge changes on one branch into another
    \item git merge <branchname>
    \item fast-forward merge
    \item merge commits
    \item conflicts can occur (merge conflicts)
\end{itemize}
\end{frame}

\begin{frame}
\frametitle{Merging branches}
\begin{center}
    \resizebox{!}{0.7\textheight}{%
        \input{figures/merge_simple_branch_tikz.tex}
    }
\end{center}
\end{frame}

\begin{frame}
\frametitle{Example: quick fix}
\begin{itemize}
    \item when working on a new feature (on a feature branch)
    \item bug found on main branch -> quick fix necessary
    \item commit changes on current branch
    \item checkout master
    \item checkout quick fix branch (git checkout -b <branchname>)
    \item make fix, test, commit
    \item merge into master, push
    \item return to work on feature branch
\end{itemize}
\end{frame}

\begin{frame}
\frametitle{Example: handling merge conflicts}
\begin{itemize}
    \item work on same area of code
    \item try to merge
    \item handle (resolve) conflict
    \item git add (to mark as resolved)
    \item git commit (as necessary)
\end{itemize}
\end{frame}

\begin{frame}
\frametitle{Cherry picking}
\begin{itemize}
    \item associated with merging
    \item only when some commits are wanted
    \item git cherry-pick <commit>
\end{itemize}
\end{frame}

\begin{frame}
\frametitle{Local and remote branches}
\end{frame}

\begin{frame}
\frametitle{Fetching from remotes}
\begin{itemize}
    \item git pull $\approx$ git fetch; git merge
    \item why fetch then merge?
    \item origin is writable but upstream isn't -> common OSS model
    \item define upstream
    \item changes are made on upstream
    \item fetch changes to local repo
    \item merge changes into local branch (e.g.\ master)
\end{itemize}
\end{frame}

\begin{frame}
\frametitle{Exercise!}
\end{frame}

\begin{frame}
\frametitle{Rewriting history}
\begin{itemize}
    \item unpublished (not yet pushed) changes can be rearranged
    \item reorder, edit, squash, reword, delete (split also possible,
        but more involved)
    \item the story on wants to tell about the development; not all
        commits need to be pushed
    \item allows one to polish the changes; improve quality
    \item linearise changes; published repo has clear, linear history
        (see e.g.\ Ovid's slide about git histories)
\end{itemize}
\end{frame}

\begin{frame}
\frametitle{Interactive rebase}
\begin{itemize}
    \item git rebase -i
    \item very handy to clean up history before pushing
\end{itemize}
\end{frame}

\begin{frame}
\frametitle{Stashing changes}
\begin{itemize}
    \item why use this?
    \item what is this?
    \item handy feature for interrupted work
    \item e.g.\ quick fix example
    \item git stash
    \item git stash show -p
    \item git stash pop (conflicts possible)
    \item git stash drop
    \item git stash drop/show {n}
\end{itemize}
\end{frame}

\begin{frame}
\frametitle{Stashing changes}
\code{git stash}
\begin{itemize}
    \item One is in the middle of a task
    \item Something unrelated needs to be fixed quickly
    \item Use \code{git stash} to save the current state (or more
        explicitly: \code{git stash save}
    \item The working copy is now "clean"
    \item One is now able to work on the new topic and check in the
        changes
    \item \code{git stash pop} reinstates the old state
    \item Now one can work on the previous task again, where one
        left off
    \item \code{git stash} works like a stack (push/pop operations)
\end{itemize}
\end{frame}

\begin{frame}
\frametitle{Stashing changes}

\begin{description}
    \item \code{git stash list} List all “stashed” states
    \item \code{git stash show} Show a stashed state
    \item \code{git stash apply} Apply a state to the current working copy
    \item \code{git stash pop} Pop the “top most” state off the stack and
        apply it onto the current working copy
\end{description}
\end{frame}

%%%%%%%%%%%%%%%%%%%%%%%%%%%%%%%%%%%%%%%%%%%%%%%%%%%%%%%%%%%%%%%%%%%%%%%%%%%%
\section{Using a version control system service}

\begin{frame}
    \frametitle{Online VCS services}
    \begin{itemize}
        \item GitHub
        \item GitLab
        \item BitBucket
        \item others?
    \end{itemize}
\end{frame}

\begin{frame}[fragile]
    \frametitle{Forking a repository}

    - create own repository of another project
    - gives ability to contribute to original project via Pull Requests
    (also called Merge Requests)
    - one can push changes to own fork because it belongs to you
    - pushing changes directly to someone else's project isn't possible
    (unless you have commit access, which isn't default situation)
    - commit changes to a non-\code{main} branch and push to upstream fork ->
    from here can create Pull Request
    - services like GitHub/GitLab detect such pushes and prompt to create a
    Pull Request
    - fork and pull request has become most common open source collaboration
    workflow (mostly due to GitHub)
    diagram:
    \begin{minted}{shell}
    |upstream repo; label: "upstream"|  -> fork -> |upstream fork; label: "origin"|
                        fetch v                            push \^ ; clone v
                                                   |local repo|
    \end{minted}

    - show screenshots of example of forking a repo on GitHub
\end{frame}

\begin{frame}
    \frametitle{Submitting a Pull Request}
    - a Pull Request is a \emph{proposal} for a code change in an upstream
    project repository.  It is merely \emph{one} way to contribute to a
    project; there are other workflows.
    - the first step is to create a fork of the project repository that
    you're contributing to (see Forking a repository)
    - in your local Git repository of the project, create a feature branch
    git checkout -b "feature\_branch\_name"
    - build your proposed feature in the feature branch
    - push feature branch to \emph{your} upstream forked repository
    - button will (usually) appear which leads to pull request dialog
    - open pull request dialog, add title (describes proposed changes very
    briefly)
    - add details of what was changed and why this is of value, was
    necessary, is a good idea in the pull request text
    - submit pull request
    - maintainer will be notified by e.g.\ GitHub automatically
    - maintainer might ask for code to be changed, or want things done
    differently. Be gracious! Try to adhere to project's development
    guidelines. Be respectful! Usually maintainers are unpaid and doing such
    projects in their spare time.
\end{frame}


\begin{frame}
    \frametitle{Pull requests in the enterprise}

The pull request model has become so popular that it has replaced other
workflows previously common in closed-source environments.  Once upon a
time, the repository would be central and everyone would be making changes
to the main/master/trunk branch all of the time.  Because everyone on
the team was working towards a common goal there wasn't much need to have a
process to say "hey, I'm adding some code to the project!" because that was
everyone's job.  And anyway, if you saw a strange change to the code you
could simply walk to a nearby office and ask your co-worker about it.  With
pull requests, there's more ceremony and friction involved in getting code
changes merged into a project.  This is both a blessing and a curse.  Using
a pull request model (automated via a centralised Git service) allows more
people to review the code changes before they land into the main branch and
hence there is more opportunity for bugs to be found and for information to
spread throughout the team.  This improves code quality and reduces the risk
of breakages once the code goes "live".  It also allows for teams to work
remotely and hence distributed across the globe, meaning we don't have to
sit in the same office to get things done anymore.  Yet, it adds more work,
friction and delay to the process of getting changes made.  Thus, fixes or
new features can be slowed down by the extra steps required to integrate
code into the main branch, which can be frustrating for developers and users
alike.

Assuming that this process is more of a blessing than a curse, it has one
technical advantage: the repository setup is a little bit easier than the
situation with Eck and Alice as described earlier.  This is because an
explicit fork for each contributor is no longer necessary.  The upstream
repository becomes a central reference point and pull requests are based off
branches pushed to the central repository rather than requiring the extra
step of pulling code from a separate upstream forked repository.  This also
makes merging code a bit easier.  Since everything revolves around the
central repository one can implement automated
\href{https://en.wikipedia.org/wiki/Continuous_integration}{continuous
integration} processes
which check and automatically merge branches after they have been reviewed
and signed off on.

\ldots the image from above (of the pull request flow) thus becomes more of a
triangle, because people are working from a central point of reference.
\end{frame}
